\section{Conclusion}
\label{sec:conclusion}

We argued that natural-language delegation creates a pre-execution responsibility-projection problem, that the projection can be made measurable by closing the responsibility space at $|J| = 12$, and that the protocol of projection $\to$ bid $\to$ clarify $\to$ select $\to$ execute $\to$ settle $\to$ reputation makes pre-execution commitment a first-class object in multi-agent LLM orchestration. On a $50$-example pilot the cross-family projection mismatch exceeds within-family stochastic variance by approximately $17\times$ (median $R = 14.4$), and all five guard hypotheses pass with $p < 0.0001$ --- the P1 measurability finding is robust. The P2 actionability question is reframed by a measurement-layer reanalysis under a redesigned 12-judge Anthropic-excluded panel and a three-condition split: projection-driven execution shows a directional disadvantage on the headline weighted-R1 settlement loss, with the cost concentrated on R1.4 (novelty mapping) and R1.7 (citation audit). A closed Stage 1 R1.7 human-anchor pilot independently surfaces a methodological boundary-condition finding (Sections~\ref{sec:results:limitations:specifiability} and \ref{sec:discussion:specifiability}): human-anchor scalability is bound by anchor specifiability and domain-expertise gating, not by rater throughput, and the form-embedded layer and the separately-read protocol-document layer are separable design moves. The convergence is direct on R1.7 (Stage 1 provides a direct boundary-condition signal on R1.7) and inferential on R1.4 (not in the Stage 1 closure). The next phase --- main-run scale-up to $n = 300$, the full four-condition Experiment 2 including task-aware routing and CLAMBER-style baselines, longitudinal reputation evaluation on real agents, and dim-specific protocol moves (retrieval-augmentation for R1.7, bounded prior-work corpus + structured comparator template for R1.4) under the specifiability constraints surfaced here --- builds directly on the closed taxonomy and the protocol formalized here.
